 \documentclass[11pt,a4paper]{article}
 \usepackage{inputenc}
 \usepackage{fontenc}
 \usepackage[ngerman]{babel}
 \usepackage{amsmath,amssymb,amsthm}
 \usepackage{geometry}
 \geometry{margin=2.5cm}
 
 % Theorem-Umgebungen
 \newtheorem{satz}{Satz}
 \newtheorem{lemma}{Lemma}
 \newtheorem{korollar}{Korollar}
 \theoremstyle{definition}
 \newtheorem{definition}{Definition}
 \newtheorem{beispiel}{Beispiel}
 \newtheorem{bemerkung}{Bemerkung}
 
 % Operatoren
 \DeclareMathOperator{\ggT}{ggT}
 \DeclareMathOperator{\rad}{rad}
 
 \begin{document}
 	
 	\section{Arithmetisches Qubit und Trapezmodell}
 	
 	Es seien $A,B,C,E\subset\mathbb{P}\setminus\{2,3\}$ paarweise disjunkte Mengen von Primzahlen. Jede Primzahl $p\mid n$ mit $p>3$ liege in genau einer der vier Familien.
 	
 	\begin{definition}[$ABC$-/$E$-Zerlegung]
 		Sei $n\in\mathbb{N}$ mit
 		\[
 		\ggT(n,6)=1.
 		\]
 		Dann schreiben wir
 		\[
 		n=n_{ABC}\,\mathcal E(n),
 		\]
 		wobei $n_{ABC}$ das Produkt aller Primfaktorpotenzen von $n$ aus den Familien
 		$A,B,C$ und $\mathcal E(n)$ das Produkt aller Primfaktorpotenzen aus der Familie $E$ bezeichnet.
 		Diese Zerlegung ist nach dem vorherigen Satz eindeutig.
 	\end{definition}
 	
 	\begin{definition}[Arithmetisches Qubit]
 		Zu einer Zahl $n\in\mathbb{N}$ mit $\ggT(n,6)=1$ definieren wir den normierten Zweizustandsvektor
 		\[
 		|\psi(n)\rangle
 		:=
 		\alpha(n)\,|E\rangle+\beta(n)\,|ABC\rangle,
 		\]
 		wobei
 		\[
 		\alpha(n):=\frac{\mathcal E(n)}{\sqrt{\mathcal E(n)^2+n_{ABC}^2}},
 		\qquad
 		\beta(n):=\frac{n_{ABC}}{\sqrt{\mathcal E(n)^2+n_{ABC}^2}}.
 		\]
 		Wir nennen $|\psi(n)\rangle$ das \emph{arithmetische Qubit} von $n$.
 	\end{definition}
 	
 	\begin{lemma}[Normierung]
 		Für jedes $n\in\mathbb{N}$ mit $\ggT(n,6)=1$ gilt
 		\[
 		\alpha(n)^2+\beta(n)^2=1.
 		\]
 	\end{lemma}
 	
 	\begin{proof}
 		Direktes Einsetzen liefert
 		\[
 		\alpha(n)^2+\beta(n)^2
 		=
 		\frac{\mathcal E(n)^2}{\mathcal E(n)^2+n_{ABC}^2}
 		+
 		\frac{n_{ABC}^2}{\mathcal E(n)^2+n_{ABC}^2}
 		=1.
 		\]
 	\end{proof}
 	
 	\begin{definition}[Unterer Raum und Trapezmodell]
 		Sei
 		\[
 		n_{ABC}=n_A\,n_B\,n_C
 		\]
 		die weitere Zerlegung des $ABC$-Anteils nach den drei Familien $A,B,C$.
 		Wir definieren den \emph{unteren Raum} von $n$ durch die drei horizontalen Komponenten
 		\[
 		n_A,\qquad n_B,\qquad n_C
 		\]
 		und die Höhe
 		\[
 		h(n):=\rad(\mathcal E(n)).
 		\]
 		Das zugehörige \emph{arithmetische Trapez} $T(n)$ sei durch die vier Punkte
 		\[
 		P_1=(0,0),\qquad
 		P_2=(n_A+n_B+n_C,0),
 		\]
 		\[
 		P_3=(n_C+n_B,h(n)),\qquad
 		P_4=(n_C,h(n))
 		\]
 		gegeben.
 	\end{definition}
 	
 	\begin{bemerkung}
 		Die Punkte $P_1,P_2,P_3,P_4$ definieren tatsächlich ein Trapez, da die Strecken
 		\[
 		\overline{P_1P_2}
 		\quad\text{und}\quad
 		\overline{P_4P_3}
 		\]
 		parallel zur $x$-Achse verlaufen. Für $n_B>0$ handelt es sich um ein echtes Trapez, für $n_B=0$ entartet es zu einem Rechteck.
 	\end{bemerkung}
 	
 	\begin{satz}[Eindeutige Zuordnung eines arithmetischen Qubits und Trapezes]
 		Jede natürliche Zahl $n\in\mathbb{N}$ mit $\ggT(n,6)=1$ bestimmt eindeutig
 		\begin{enumerate}
 			\item ihr arithmetisches Qubit $|\psi(n)\rangle$,
 			\item ihr arithmetisches Trapez $T(n)$ im unteren Raum.
 		\end{enumerate}
 	\end{satz}
 	
 	\begin{proof}
 		Die Größen $n_{ABC}$ und $\mathcal E(n)$ sind durch die kanonische $ABC$-/$E$-Zerlegung eindeutig bestimmt. Damit sind auch
 		\[
 		\alpha(n),\qquad \beta(n)
 		\]
 		eindeutig bestimmt, also das arithmetische Qubit $|\psi(n)\rangle$.
 		
 		Ebenso ist die Zerlegung
 		\[
 		n_{ABC}=n_A\,n_B\,n_C
 		\]
 		nach den drei Familien eindeutig. Daraus folgen die Koordinaten
 		\[
 		P_1,\ P_2,\ P_3,\ P_4
 		\]
 		des Trapezes eindeutig, da auch
 		\[
 		h(n)=\rad(\mathcal E(n))
 		\]
 		eindeutig bestimmt ist. Also ist $T(n)$ eindeutig.
 	\end{proof}
 	
 	\begin{korollar}
 		Jede natürliche Zahl $n$ mit $\ggT(n,6)=1$ besitzt eine eindeutige zweifache Repräsentation:
 		\begin{enumerate}
 			\item algebraisch durch ihre kanonische Zerlegung
 			\[
 			n=n_{ABC}\,\mathcal E(n),
 			\]
 			\item geometrisch-qubitartig durch das Paar
 			\[
 			\bigl(|\psi(n)\rangle,\;T(n)\bigr).
 			\]
 		\end{enumerate}
 	\end{korollar}
 	
 	\begin{bemerkung}
 		Diese Konstruktion identifiziert eine Zahl nicht mit einem physikalischen Qubit, sondern mit
 		einem \emph{arithmetisch definierten Zwei-Zustands-Objekt}. Ebenso ist das Trapezmodell eine
 		kanonische geometrische Repräsentation der Zerlegung, nicht bereits eine zusätzliche
 		zahlentheoretische Invarianzbehauptung.
 	\end{bemerkung}
 	
 	\begin{beispiel}
 		Für
 		\[
 		n=5\cdot 7\cdot 11\cdot 13^2
 		\]
 		gilt
 		\[
 		n_A=5,\qquad n_B=7,\qquad n_C=11,
 		\qquad
 		\mathcal E(n)=13^2,
 		\]
 		also
 		\[
 		n_{ABC}=5\cdot 7\cdot 11=385.
 		\]
 		Damit ist
 		\[
 		|\psi(n)\rangle
 		=
 		\frac{169}{\sqrt{169^2+385^2}}\,|E\rangle
 		+
 		\frac{385}{\sqrt{169^2+385^2}}\,|ABC\rangle.
 		\]
 		Ferner ist
 		\[
 		h(n)=\rad(13^2)=13,
 		\]
 		und das zugehörige Trapez ist gegeben durch
 		\[
 		P_1=(0,0),\quad
 		P_2=(23,0),\quad
 		P_3=(18,13),\quad
 		P_4=(11,13).
 		\]
 	\end{beispiel}
 	
 \end{document}